\begin{table}[!htb]
\centering
\setlength{\tabcolsep}{3mm}
\begin{sc}
\begin{tabular}{c|cc}
\toprule
Model & FLOPs & OPs \\
\midrule
Transformer & 80.04M & - \\
GRU & 37.49M & - \\
RNN & 35.36M & - \\
RDDLGN & - & 1.53M \\
\bottomrule
\end{tabular}
\end{sc}
\caption{Computational complexity comparison across different model architectures (excluding embedding operations). 
FLOPs (Floating Point Operations) are calculated for a single forward pass with batch size 1, 
source sequence length 16, target sequence length 16, and vocabulary size 16,000. 
Embedding lookup operations are excluded from all FLOP calculations to focus on core architectural differences. 
For RNN and GRU models, FLOPs include recurrent computations (matrix multiplications, element-wise operations, 
and activations) and output projections. For Transformer models, FLOPs include positional encodings, 
multi-head attention mechanisms (Q/K/V projections, attention scores, softmax operations), feed-forward networks, 
layer normalizations, and output projections across all encoder and decoder layers. OPs (Operations) for RDDLGN 
represent the sum of all layer sizes from the model configuration, indicating the total number of logical 
operations in the differentiable logic network.}
\label{tab:computational_complexity}
\end{table}